贪心证明深度题卡 按题型拆成章内大节，但每个大节都先从 LeetCode 案例切入，再进入分流、案例矩阵、案例推导和实验复盘。这样读者看到的是从具体题推到规律，而不是先读抽象方法再猜它能用在哪。

\topic{贪心与交换证明}
这一节先看 LeetCode 45 跳跃游戏 II、LeetCode 55 跳跃游戏、LeetCode 122 买卖股票 II 这几道 LeetCode 题，再整理同一题型的分流和推导。阅读顺序不是先背原理，而是先看案例的暴力瓶颈，再把重复出现的观察抽象成方法。

\subsection{原理和适用条件}

以 LeetCode 45 跳跃游戏 II、LeetCode 55 跳跃游戏、LeetCode 122 买卖股票 II 为主线：LeetCode 45 跳跃游戏 II：模型是“每一次跳跃可以看成 BFS 的一层，当前层内所有位置共享同一次跳数。”，优化抓手是“扫描当前层时维护下一层最远边界，到达当前层末尾才增加跳数。”；LeetCode 55 跳跃游戏：模型是“扫描过程中只关心当前能到达的最远位置。”，优化抓手是“如果下标超过最远可达，任何之前路径都不能到达它；否则用它扩展边界。”；LeetCode 122 买卖股票 II：模型是“允许多次交易且不能同时持有，所有正收益差都可以收集。”，优化抓手是“把长上涨段拆成每天买卖收益相同，因此累加正差值最优。”。这些案例共同推出的规律是：先把选择空间完整枚举，再寻找可以交换到最优解前面的局部选择。贪心必须能证明当前选择不会让后续资源变差，而不是只看排序后代码短。 方法名只是复盘后的标签，真正要掌握的是从题面约束、暴力失败、状态设计到边界反例的推导链。

\begin{principlebox}{图示：从 LeetCode 案例到 贪心与交换证明推导链}
\begin{tabular}{@{}p{0.18\linewidth}p{0.74\linewidth}@{}}
暴力基线 & 枚举候选、完整模拟或逐个检查，先保证覆盖全部可能答案。\\
瓶颈定位 & 慢点是选择空间爆炸。若每一步都有很多候选，而某个排序规则或边界规则能安全丢掉大部分候选，就可能存在贪心。\\
结构选择 & 贪心需要找到可交换的局部最优：最早结束、最小右端、当前最远覆盖、当前最大收益或最小代价。排序只是手段，不是证明。\\
不变量 & 贪心不变量通常是当前方案在相同选择次数下不落后于任何其他方案，或者当前保留的资源最多、结束最早、右边界最小。\\
代码落地 & 实现时先把比较器写成明确的业务顺序。区间题注意起点相同时终点升序还是降序；覆盖题注意闭区间和半开区间的等号。\\
\end{tabular}
\end{principlebox}

\subsection{LeetCode 案例分流}

\begin{principlebox}{从 LeetCode 案例分流：贪心与交换证明}
\textbf{最早结束和区间选择。}
代表案例：LeetCode 435 无重叠区间、LeetCode 452 引爆气球、LeetCode 646 最长数对链。先看这些题的题面条件：选择一个对象会占用时间或空间区间。由此推出的处理方式是：按结束位置排序，优先保留结束最早的候选。风险点：必须用交换论证，不是所有按端点排序都正确。

\textbf{最远覆盖。}
代表案例：LeetCode 55 跳跃游戏、LeetCode 45 跳跃游戏 II、LeetCode 763 划分字母区间。先看这些题的题面条件：当前选择决定下一段可覆盖范围。由此推出的处理方式是：扫描时维护当前可达最远边界或当前片段右端。风险点：当前层边界和下一层最远边界不能混。

\textbf{局部约束两遍扫描。}
代表案例：LeetCode 135 分发糖果、LeetCode 42 接雨水前后缀思路。先看这些题的题面条件：相邻约束可从左右两个方向分别满足。由此推出的处理方式是：左到右满足一侧约束，右到左满足另一侧，最后合并。风险点：只扫一遍往往不能同时满足两侧关系。

\textbf{资源反悔。}
代表案例：LeetCode 630 课程表 III、LeetCode 871 最低加油次数。先看这些题的题面条件：当前选了某些对象后，未来遇到更优对象可能需要替换。由此推出的处理方式是：先接受候选，用堆或排序在超约束时丢掉最差选择。风险点：反悔贪心要证明丢掉的候选不会让结果变差。

\end{principlebox}

\subsection{案例矩阵}

本节案例覆盖 LeetCode 45 跳跃游戏 II、LeetCode 55 跳跃游戏、LeetCode 122 买卖股票 II、LeetCode 134 加油站、LeetCode 135 分发糖果、LeetCode 169 多数元素、LeetCode 316 去除重复字母、LeetCode 406 根据身高重建队列、LeetCode 435 无重叠区间、LeetCode 452 用最少数量的箭引爆气球、LeetCode 630 课程表 III、LeetCode 646 最长数对链、LeetCode 678 有效的括号字符串、LeetCode 763 划分字母区间、LeetCode 860 柠檬水找零。阅读时不要按题号背答案，而要先判断题面落在哪个分流场景：查询目标是什么，输入约束给了什么单调性或等价关系，暴力慢在哪里，优化结构保存了什么状态。

\begin{principlebox}{案例矩阵}
\textbf{LeetCode 45 跳跃游戏 II。}
每一次跳跃可以看成 BFS 的一层，当前层内所有位置共享同一次跳数。单点风险：数组长度一、刚好到达、存在多个可选跳点、不要在每个位置都加跳数。

\textbf{LeetCode 55 跳跃游戏。}
扫描过程中只关心当前能到达的最远位置。单点风险：第一个元素为零、数组长度一、恰好到达最后、远超最后。

\textbf{LeetCode 122 买卖股票 II。}
允许多次交易且不能同时持有，所有正收益差都可以收集。单点风险：下降段、平台价格、只有一天、手续费不存在。

\textbf{LeetCode 134 加油站。}
绕行一圈是否可行取决于总油量差和局部剩余油量。单点风险：总和刚好为零、多个可行起点、某站油量为零、消耗大于补给。

\textbf{LeetCode 135 分发糖果。}
相邻评分约束需要同时满足左邻和右邻两个方向。单点风险：评分相等、严格递增、严格递减、山峰和谷底。

\textbf{LeetCode 169 多数元素。}
超过一半的元素可以和其他元素两两抵消后仍然剩下。单点风险：题目是否保证存在多数、计数归零、全相同。

\textbf{LeetCode 316 去除重复字母。}
每个字符必须保留一次，同时结果字典序尽量小。单点风险：字符只出现一次、重复字符、弹出后 visited 恢复、字典序比较。

\textbf{LeetCode 406 根据身高重建队列。}
高个子的位置先确定后，不会被矮个子影响其前方高个数量。单点风险：相同身高、k 为零、插入到末尾、结果稳定性。

\textbf{LeetCode 435 无重叠区间。}
保留最多不重叠区间等价于删除最少区间。单点风险：端点相接是否重叠、完全覆盖、相同终点、空列表。

\textbf{LeetCode 452 用最少数量的箭引爆气球。}
一支箭的位置可以覆盖所有包含该位置的区间。单点风险：端点相同、负坐标、区间包含、闭区间边界。

\textbf{LeetCode 630 课程表 III。}
每门课有持续时间和截止日，已选课程总时长不能超过当前截止日。单点风险：课程无法单独完成、截止日相同、替换已选课程、空列表。

\textbf{LeetCode 646 最长数对链。}
每个数对像一个区间，选择链要求前一个右端小于后一个左端。单点风险：端点相等不允许连接、完全覆盖、负数端点、空数组。

\textbf{LeetCode 678 有效的括号字符串。}
星号可以作为左括号、右括号或空，使未匹配左括号数量变成一个可能区间。单点风险：上界变负、最终下界不为零、连续星号、前缀右括号过多。

\textbf{LeetCode 763 划分字母区间。}
每个片段必须覆盖其中所有字符的最后出现位置。单点风险：单字符、所有字符相同、字符多次跨段。

\textbf{LeetCode 860 柠檬水找零。}
找零资源只有 5 元和 10 元，5 元比 10 元更灵活。单点风险：第一张不是 5、连续 10、连续 20、只有 5、5 元数量不足。

\end{principlebox}

\subsection{共性落点}

\begin{principlebox}{本组共性落点}
\textbf{慢版本。} 暴力做法枚举所有选择顺序或所有子集，找出最优方案。贪心题的难点不在写循环，而在证明局部选择不会破坏全局最优。
\textbf{瓶颈。} 慢点是选择空间爆炸。若每一步都有很多候选，而某个排序规则或边界规则能安全丢掉大部分候选，就可能存在贪心。
\textbf{状态字段。} 处理顺序、当前已选方案、剩余资源或最远边界、答案计数；若使用排序，要说明排序键；每次局部选择后要能说明当前状态不落后。
\textbf{不变量。} 贪心不变量通常是当前方案在相同选择次数下不落后于任何其他方案，或者当前保留的资源最多、结束最早、右边界最小。
\textbf{实现检查。} 把局部规则写成业务含义；若需要排序，就处理相同关键字；循环中只维护证明需要的状态，不把局部选择和答案更新混在一起。
\textbf{C++ 落地。} 实现时先把比较器写成明确的业务顺序。区间题注意起点相同时终点升序还是降序；覆盖题注意闭区间和半开区间的等号。
\textbf{证明抓手。} 证明从交换或领先性开始：任意最优解都能替换成当前贪心选择而不变差，或贪心状态始终不落后。
\end{principlebox}

\subsection{案例推导}

\noindent\textbf{LeetCode 45 跳跃游戏 II。}

\textbf{题目说明。} LeetCode 45 跳跃游戏 II 的题面入口要先还原成具体任务：每一次跳跃可以看成 BFS 的一层，当前层内所有位置共享同一次跳数。

\textbf{样例入口。} 拿 [2, 3, 1, 1, 4] 手算，起点这一层能到下标 1 和 2；扫描这一层时，下一层最远可到 4，走完当前层末尾才把步数加一。

\textbf{暴力基线。} 暴力 BFS 可以从每个当前下标扩展所有可跳位置，按层找第一次到终点的步数。

\textbf{暴力过程。} 以 [2, 3, 1, 1, 4] 为例，第一层是从 0 能到的 1、2；扫描这一层时，下一层最远可到 4，所以答案是 2。

\textbf{重复工作。} 题面模型：每一次跳跃可以看成 BFS 的一层，当前层内所有位置共享同一次跳数。暴力版的慢点要落到本题：浪费在显式保存一层内所有节点。数组下标的可达层是连续区间，只需要当前层右边界和下一层最远边界。后面的优化只消掉这类重复工作，不能改变答案集合。

\textbf{优化条件。} 本题的关键观察是：扫描当前层时维护下一层最远边界，到达当前层末尾才增加跳数。这句话必须能翻译成可维护状态；如果题目条件改变后观察不再成立，就不能继续套原来的写法。

\textbf{相邻状态。} 规律是把可达区间看成 BFS 层，当前层结束时切到下一层；不能在每个位置都加跳数。把这个特例继续拆开看：先构造一个非贪心选择，再比较两者留下的资源、右边界或剩余时间。只有能把非贪心第一步交换成贪心第一步且不变差，局部选择才是规律。

\textbf{状态复用。} current\_end 表示当前跳数能覆盖的最右位置，farthest 表示下一跳能覆盖的最远位置。扫描到 current\_end 时，步数加一并切换到下一层。手算时只改被新输入影响的那一小部分，而不是回头重算完整候选。

\textbf{指针和字段。} 状态字段要能说清：i 是扫描位置，jumps 是已完成层数，current\_end 是本层末尾，farthest 是下一层末尾。手算时按这个协议跟踪：i=0 时 farthest=2，走到 current\_end=0 后 jumps=1、current\_end=2；扫描 i=1 时 farthest 更新到 4，走到 i=2 后 jumps=2。

\textbf{代码落点。} 关键是循环通常只扫到 n-2，因为到达最后一个位置不需要再跳。不能在每个位置都 jumps++，只能在层结束时加。关键代码行必须能逐句翻译成“旧状态怎样变成新状态”，否则就只是模板。

\textbf{正确性。} 所有当前跳数可达的位置构成一个连续区间；在这个区间内选择能让下一层覆盖最远的位置，不会比保存全部 BFS 节点差。

\textbf{边界实验。} 先用数组长度一、刚好到达、存在多个可选跳点、不要在每个位置都加跳数做反例检查。实验时除了数组长度一、刚好到达、存在多个可选跳点、不要在每个位置都加跳数，还要故意替换成一个看似合理但错误的局部规则，构造反例。能解释错误贪心为何失败，才算理解正确贪心。

\textbf{复杂度变化。} 显式 BFS 分支多；区间层扫描 O(n) 时间、O(1) 空间。

\textbf{迁移追问。} 如果题目改成“若问能否到达，只需维护全局最远可达位置”，先判断原状态是否仍然覆盖新答案集合，再决定是微调边界、改 value 语义，还是换算法。

\noindent\textbf{LeetCode 55 跳跃游戏。}

\textbf{题目说明。} LeetCode 55 跳跃游戏 的题面入口要先还原成具体任务：扫描过程中只关心当前能到达的最远位置。

\textbf{样例入口。} 拿 [2, 3, 1, 1, 4] 手算，下标 0 能覆盖到 2；走到下标 1 时最远可达更新到 4，已经覆盖终点。再拿 [0, 1]，扫描到下标 1 时已经超过最远可达，失败。

\textbf{暴力基线。} 暴力把每个下标能跳到的位置都当成分支搜索，尝试是否存在一条路径能到终点。

\textbf{暴力过程。} 以 [2, 3, 1, 1, 4] 为例，从 0 可以跳到 1 或 2；若到 1，最远能继续覆盖到 4。以 [0, 1] 为例，从 0 一步也走不出去。

\textbf{重复工作。} 题面模型：扫描过程中只关心当前能到达的最远位置。暴力版的慢点要落到本题：浪费在枚举具体路径。判断能否到达终点时，不关心走的是哪条路径，只关心目前所有可达位置能覆盖到的最远下标。后面的优化只消掉这类重复工作，不能改变答案集合。

\textbf{优化条件。} 本题的关键观察是：如果下标超过最远可达，任何之前路径都不能到达它；否则用它扩展边界。这句话必须能翻译成可维护状态；如果题目条件改变后观察不再成立，就不能继续套原来的写法。

\textbf{相邻状态。} 规律是只维护当前能到达的最远边界，任何超过边界的位置都不可能由前面路径到达。把这个特例继续拆开看：先构造一个非贪心选择，再比较两者留下的资源、右边界或剩余时间。只有能把非贪心第一步交换成贪心第一步且不变差，局部选择才是规律。

\textbf{状态复用。} 扫描下标 i，维护 farthest。若 i>farthest，说明当前下标不可达，直接失败；否则用 i+nums[i] 更新 farthest。手算时只改被新输入影响的那一小部分，而不是回头重算完整候选。

\textbf{指针和字段。} 状态字段要能说清：i 是当前检查位置，farthest 是所有已可达位置能跳到的最远边界，n-1 是目标边界。手算时按这个协议跟踪：样例中 i=0 后 farthest=2；i=1 可达并把 farthest 更新为 4，已经覆盖终点。

\textbf{代码落点。} 关键顺序是先判断当前位置是否超过 farthest；若没超过，再用 i+nums[i] 扩展最远可达位置。数组长度为 1 时起点就是终点。关键代码行必须能逐句翻译成“旧状态怎样变成新状态”，否则就只是模板。

\textbf{正确性。} 在扫描到 i 前，farthest 汇总了所有可达下标的最远覆盖；若 i 超过它，任何旧路径都到不了 i，更不可能通过 i 之后的位置到终点。

\textbf{边界实验。} 先用第一个元素为零、数组长度一、恰好到达最后、远超最后做反例检查。实验时除了第一个元素为零、数组长度一、恰好到达最后、远超最后，还要故意替换成一个看似合理但错误的局部规则，构造反例。能解释错误贪心为何失败，才算理解正确贪心。

\textbf{复杂度变化。} 路径搜索最坏指数级；贪心边界扫描 O(n) 时间、O(1) 空间。

\textbf{迁移追问。} 如果题目改成“求最少跳数时把可达区间看成 BFS 层”，先判断原状态是否仍然覆盖新答案集合，再决定是微调边界、改 value 语义，还是换算法。

\noindent\textbf{LeetCode 122 买卖股票 II。}

\textbf{题目说明。} LeetCode 122 买卖股票 II 的题面入口要先还原成具体任务：允许多次交易且不能同时持有，所有正收益差都可以收集。

\textbf{样例入口。} 拿 [7, 1, 5, 3, 6, 4] 手算，1 到 5 的利润 4，加上 3 到 6 的利润 3，总利润 7。把 1 到 6 的上涨段拆成每天正差值，收益相同。

\textbf{暴力基线。} 慢版本枚举选择顺序或用 DP 保留多个可能方案，先确认最优值存在。随后才检查局部选择是否能通过交换或领先性固定下来。放到本题就是：先按“允许多次交易且不能同时持有，所有正收益差都可以收集”完整试慢版本；样例里已经能看到“规律是无限次交易且无手续费时，所有正差值都可收集；下降段和平台不贡献收益”，所以优化前必须先能说清慢版本枚举了哪些候选。

\textbf{暴力过程。} 拿 [7, 1, 5, 3, 6, 4] 手算，1 到 5 的利润 4，加上 3 到 6 的利润 3，总利润 7。把 1 到 6 的上涨段拆成每天正差值，收益相同。

\textbf{重复工作。} 题面模型：允许多次交易且不能同时持有，所有正收益差都可以收集。暴力版的慢点要落到本题：慢版本会围绕“允许多次交易且不能同时持有，所有正收益差都可以收集”枚举选择顺序。样例规律是“规律是无限次交易且无手续费时，所有正差值都可收集；下降段和平台不贡献收益”，说明存在一个局部选择或边界状态可以安全保留；不能证明这个选择不变差时，就不能把它叫贪心。后面的优化只消掉这类重复工作，不能改变答案集合。

\textbf{优化条件。} 本题的关键观察是：把长上涨段拆成每天买卖收益相同，因此累加正差值最优。这句话必须能翻译成可维护状态；如果题目条件改变后观察不再成立，就不能继续套原来的写法。

\textbf{相邻状态。} 规律是无限次交易且无手续费时，所有正差值都可收集；下降段和平台不贡献收益。把这个特例继续拆开看：先构造一个非贪心选择，再比较两者留下的资源、右边界或剩余时间。只有能把非贪心第一步交换成贪心第一步且不变差，局部选择才是规律。

\textbf{状态复用。} 把指数级选择压成一个可证明安全的局部选择；状态通常是当前最远边界、最早结束时间、最小代价或剩余资源。本题落点是：规律是无限次交易且无手续费时，所有正差值都可收集；下降段和平台不贡献收益。手算时只改被新输入影响的那一小部分，而不是回头重算完整候选。

\textbf{指针和字段。} 状态字段要能说清：处理顺序、当前已选方案、剩余资源或最远边界、答案计数；若使用排序，要说明排序键；每次局部选择后要能说明当前状态不落后。手算时按这个协议跟踪：拿 [7, 1, 5, 3, 6, 4] 手算，1 到 5 的利润 4，加上 3 到 6 的利润 3，总利润 7。把 1 到 6 的上涨段拆成每天正差值，收益相同。然后按协议复查：手算时同时写一个非贪心选择，与贪心选择比较剩余资源。若无法说明贪心后剩余资源不差，就不能直接下结论。

\textbf{代码落点。} 把局部规则写成业务含义；若需要排序，就处理相同关键字；循环中只维护证明需要的状态，不把局部选择和答案更新混在一起。关键代码行必须能逐句翻译成“旧状态怎样变成新状态”，否则就只是模板。

\textbf{正确性。} 证明从交换或领先性开始：任意最优解都能替换成当前贪心选择而不变差，或贪心状态始终不落后。

\textbf{边界实验。} 先用下降段、平台价格、只有一天、手续费不存在做反例检查。实验时除了下降段、平台价格、只有一天、手续费不存在，还要故意替换成一个看似合理但错误的局部规则，构造反例。能解释错误贪心为何失败，才算理解正确贪心。

\textbf{复杂度变化。} 暴力枚举选择顺序可能指数级；贪心通常先排序，再线性扫描或用堆维护少量候选。

\textbf{迁移追问。} 如果题目改成“有手续费时不能简单累加正差，需要状态机或调整贪心阈值”，先判断原状态是否仍然覆盖新答案集合，再决定是微调边界、改 value 语义，还是换算法。

\noindent\textbf{LeetCode 134 加油站。}

\textbf{题目说明。} LeetCode 134 加油站 的题面入口要先还原成具体任务：绕行一圈是否可行取决于总油量差和局部剩余油量。

\textbf{样例入口。} 拿 gas=[1, 2, 3, 4, 5]、cost=[3, 4, 5, 1, 2] 手算，从 0 出发中途油量为负，0 到失败点之间任何站都不能作为起点；从 3 出发可以绕完。

\textbf{暴力基线。} 暴力把每个加油站都当起点，模拟绕一圈，若油量中途为负就换下一个起点。

\textbf{暴力过程。} gas=[1, 2, 3, 4, 5]、cost=[3, 4, 5, 1, 2] 中，从 0 出发会在中途油量为负；从 3 出发可以完成一圈。

\textbf{重复工作。} 题面模型：绕行一圈是否可行取决于总油量差和局部剩余油量。暴力版的慢点要落到本题：浪费在失败起点之后仍逐个重试中间站。若从 start 到 i 的累计油量已经为负，那么 start..i 中任何站都不可能作为可行起点。后面的优化只消掉这类重复工作，不能改变答案集合。

\textbf{优化条件。} 本题的关键观察是：总和为负必不可行；扫描时若当前剩余为负，起点只能放到下一站。这句话必须能翻译成可维护状态；如果题目条件改变后观察不再成立，就不能继续套原来的写法。

\textbf{相邻状态。} 规律是总油量差为负则无解，局部剩余为负时起点跳到下一站。把这个特例继续拆开看：先构造一个非贪心选择，再比较两者留下的资源、右边界或剩余时间。只有能把非贪心第一步交换成贪心第一步且不变差，局部选择才是规律。

\textbf{状态复用。} 扫描一次，total 记录总油量差，tank 记录当前候选起点到当前位置的剩余油量，start 是当前候选起点。手算时只改被新输入影响的那一小部分，而不是回头重算完整候选。

\textbf{指针和字段。} 状态字段要能说清：diff=gas[i]-cost[i]，tank 是局部剩余，total 是全局可行性，start 是下一次尝试起点。手算时按这个协议跟踪：样例前几站累计 tank 变负后，把 start 移到下一站并清零 tank；最后 total>=0，返回最终 start=3。

\textbf{代码落点。} 关键是 total 全程累加，tank<0 时 start=i+1、tank=0。循环结束后 total<0 返回 -1，否则返回 start。关键代码行必须能逐句翻译成“旧状态怎样变成新状态”，否则就只是模板。

\textbf{正确性。} tank<0 说明从当前 start 到 i 的总补给不足；其中任意中间站作为起点，少了前面一段补给，只会更差，故可整体跳过。

\textbf{边界实验。} 先用总和刚好为零、多个可行起点、某站油量为零、消耗大于补给做反例检查。实验时除了总和刚好为零、多个可行起点、某站油量为零、消耗大于补给，还要故意替换成一个看似合理但错误的局部规则，构造反例。能解释错误贪心为何失败，才算理解正确贪心。

\textbf{复杂度变化。} 暴力 O(\ensuremath{n^{2}})，一次扫描 O(n) 时间、O(1) 空间。

\textbf{迁移追问。} 如果题目改成“它的领先性证明是前一段任意站点都不能作为起点”，先判断原状态是否仍然覆盖新答案集合，再决定是微调边界、改 value 语义，还是换算法。

\noindent\textbf{LeetCode 135 分发糖果。}

\textbf{题目说明。} LeetCode 135 分发糖果 的题面入口要先还原成具体任务：相邻评分约束需要同时满足左邻和右邻两个方向。

\textbf{样例入口。} 拿 ratings=[1, 0, 2] 手算，中间孩子评分最低，左右都要比它多糖，结果 [2, 1, 2]。只从左到右会漏掉右侧约束。

\textbf{暴力基线。} 慢版本枚举选择顺序或用 DP 保留多个可能方案，先确认最优值存在。随后才检查局部选择是否能通过交换或领先性固定下来。放到本题就是：先按“相邻评分约束需要同时满足左邻和右邻两个方向”完整试慢版本；样例里已经能看到“规律是左到右满足左邻递增，右到左满足右邻递增，最终取两侧要求最大值”，所以优化前必须先能说清慢版本枚举了哪些候选。

\textbf{暴力过程。} 拿 ratings=[1, 0, 2] 手算，中间孩子评分最低，左右都要比它多糖，结果 [2, 1, 2]。只从左到右会漏掉右侧约束。

\textbf{重复工作。} 题面模型：相邻评分约束需要同时满足左邻和右邻两个方向。暴力版的慢点要落到本题：慢版本会围绕“相邻评分约束需要同时满足左邻和右邻两个方向”枚举选择顺序。样例规律是“规律是左到右满足左邻递增，右到左满足右邻递增，最终取两侧要求最大值”，说明存在一个局部选择或边界状态可以安全保留；不能证明这个选择不变差时，就不能把它叫贪心。后面的优化只消掉这类重复工作，不能改变答案集合。

\textbf{优化条件。} 本题的关键观察是：左到右满足递增关系，右到左满足反向递增关系，最后取两侧要求最大值。这句话必须能翻译成可维护状态；如果题目条件改变后观察不再成立，就不能继续套原来的写法。

\textbf{相邻状态。} 规律是左到右满足左邻递增，右到左满足右邻递增，最终取两侧要求最大值。把这个特例继续拆开看：先构造一个非贪心选择，再比较两者留下的资源、右边界或剩余时间。只有能把非贪心第一步交换成贪心第一步且不变差，局部选择才是规律。

\textbf{状态复用。} 把指数级选择压成一个可证明安全的局部选择；状态通常是当前最远边界、最早结束时间、最小代价或剩余资源。本题落点是：规律是左到右满足左邻递增，右到左满足右邻递增，最终取两侧要求最大值。手算时只改被新输入影响的那一小部分，而不是回头重算完整候选。

\textbf{指针和字段。} 状态字段要能说清：处理顺序、当前已选方案、剩余资源或最远边界、答案计数；若使用排序，要说明排序键；每次局部选择后要能说明当前状态不落后。手算时按这个协议跟踪：拿 ratings=[1, 0, 2] 手算，中间孩子评分最低，左右都要比它多糖，结果 [2, 1, 2]。只从左到右会漏掉右侧约束。然后按协议复查：手算时同时写一个非贪心选择，与贪心选择比较剩余资源。若无法说明贪心后剩余资源不差，就不能直接下结论。

\textbf{代码落点。} 把局部规则写成业务含义；若需要排序，就处理相同关键字；循环中只维护证明需要的状态，不把局部选择和答案更新混在一起。关键代码行必须能逐句翻译成“旧状态怎样变成新状态”，否则就只是模板。

\textbf{正确性。} 证明从交换或领先性开始：任意最优解都能替换成当前贪心选择而不变差，或贪心状态始终不落后。

\textbf{边界实验。} 先用评分相等、严格递增、严格递减、山峰和谷底做反例检查。实验时除了评分相等、严格递增、严格递减、山峰和谷底，还要故意替换成一个看似合理但错误的局部规则，构造反例。能解释错误贪心为何失败，才算理解正确贪心。

\textbf{复杂度变化。} 暴力枚举选择顺序可能指数级；贪心通常先排序，再线性扫描或用堆维护少量候选。

\textbf{迁移追问。} 如果题目改成“这是局部约束两遍扫描，不能只用一个方向贪心”，先判断原状态是否仍然覆盖新答案集合，再决定是微调边界、改 value 语义，还是换算法。

\noindent\textbf{LeetCode 169 多数元素。}

\textbf{题目说明。} LeetCode 169 多数元素 的题面入口要先还原成具体任务：超过一半的元素可以和其他元素两两抵消后仍然剩下。

\textbf{样例入口。} 拿 [2, 2, 1, 1, 1, 2, 2] 手算，候选和不同元素互相抵消，最后剩下的候选是 2。

\textbf{暴力基线。} 慢版本枚举选择顺序或用 DP 保留多个可能方案，先确认最优值存在。随后才检查局部选择是否能通过交换或领先性固定下来。放到本题就是：先按“超过一半的元素可以和其他元素两两抵消后仍然剩下”完整试慢版本；样例里已经能看到“规律是多数元素超过一半，和其他元素两两抵消后仍会剩下；题目不保证存在时需要二次计数验证”，所以优化前必须先能说清慢版本枚举了哪些候选。

\textbf{暴力过程。} 拿 [2, 2, 1, 1, 1, 2, 2] 手算，候选和不同元素互相抵消，最后剩下的候选是 2。

\textbf{重复工作。} 题面模型：超过一半的元素可以和其他元素两两抵消后仍然剩下。暴力版的慢点要落到本题：慢版本会围绕“超过一半的元素可以和其他元素两两抵消后仍然剩下”枚举选择顺序。样例规律是“规律是多数元素超过一半，和其他元素两两抵消后仍会剩下；题目不保证存在时需要二次计数验证”，说明存在一个局部选择或边界状态可以安全保留；不能证明这个选择不变差时，就不能把它叫贪心。后面的优化只消掉这类重复工作，不能改变答案集合。

\textbf{优化条件。} 本题的关键观察是：Boyer-Moore 投票维护候选和计数。这句话必须能翻译成可维护状态；如果题目条件改变后观察不再成立，就不能继续套原来的写法。

\textbf{相邻状态。} 规律是多数元素超过一半，和其他元素两两抵消后仍会剩下；题目不保证存在时需要二次计数验证。把这个特例继续拆开看：先构造一个非贪心选择，再比较两者留下的资源、右边界或剩余时间。只有能把非贪心第一步交换成贪心第一步且不变差，局部选择才是规律。

\textbf{状态复用。} 把指数级选择压成一个可证明安全的局部选择；状态通常是当前最远边界、最早结束时间、最小代价或剩余资源。本题落点是：规律是多数元素超过一半，和其他元素两两抵消后仍会剩下；题目不保证存在时需要二次计数验证。手算时只改被新输入影响的那一小部分，而不是回头重算完整候选。

\textbf{指针和字段。} 状态字段要能说清：处理顺序、当前已选方案、剩余资源或最远边界、答案计数；若使用排序，要说明排序键；每次局部选择后要能说明当前状态不落后。手算时按这个协议跟踪：拿 [2, 2, 1, 1, 1, 2, 2] 手算，候选和不同元素互相抵消，最后剩下的候选是 2。然后按协议复查：手算时同时写一个非贪心选择，与贪心选择比较剩余资源。若无法说明贪心后剩余资源不差，就不能直接下结论。

\textbf{代码落点。} 把局部规则写成业务含义；若需要排序，就处理相同关键字；循环中只维护证明需要的状态，不把局部选择和答案更新混在一起。关键代码行必须能逐句翻译成“旧状态怎样变成新状态”，否则就只是模板。

\textbf{正确性。} 证明从交换或领先性开始：任意最优解都能替换成当前贪心选择而不变差，或贪心状态始终不落后。

\textbf{边界实验。} 先用题目是否保证存在多数、计数归零、全相同做反例检查。实验时除了题目是否保证存在多数、计数归零、全相同，还要故意替换成一个看似合理但错误的局部规则，构造反例。能解释错误贪心为何失败，才算理解正确贪心。

\textbf{复杂度变化。} 暴力枚举选择顺序可能指数级；贪心通常先排序，再线性扫描或用堆维护少量候选。

\textbf{迁移追问。} 如果题目改成“若找超过三分之一的元素，需要维护两个候选”，先判断原状态是否仍然覆盖新答案集合，再决定是微调边界、改 value 语义，还是换算法。

\noindent\textbf{LeetCode 316 去除重复字母。}

\textbf{题目说明。} LeetCode 316 去除重复字母 的题面入口要先还原成具体任务：每个字符必须保留一次，同时结果字典序尽量小。

\textbf{样例入口。} 拿 cbacdcbc 手算，遇到 b 时，前面的 c 后面还会出现且 c>b，所以可以弹出 c；但遇到只剩最后一次的字符就不能弹。

\textbf{暴力基线。} 慢版本枚举选择顺序或用 DP 保留多个可能方案，先确认最优值存在。随后才检查局部选择是否能通过交换或领先性固定下来。放到本题就是：先按“每个字符必须保留一次，同时结果字典序尽量小”完整试慢版本；样例里已经能看到“规律是单调栈维护结果字典序，弹出条件必须同时满足栈顶更大且后面还会出现”，所以优化前必须先能说清慢版本枚举了哪些候选。

\textbf{暴力过程。} 拿 cbacdcbc 手算，遇到 b 时，前面的 c 后面还会出现且 c>b，所以可以弹出 c；但遇到只剩最后一次的字符就不能弹。

\textbf{重复工作。} 题面模型：每个字符必须保留一次，同时结果字典序尽量小。暴力版的慢点要落到本题：慢版本会围绕“每个字符必须保留一次，同时结果字典序尽量小”枚举选择顺序。样例规律是“规律是单调栈维护结果字典序，弹出条件必须同时满足栈顶更大且后面还会出现”，说明存在一个局部选择或边界状态可以安全保留；不能证明这个选择不变差时，就不能把它叫贪心。后面的优化只消掉这类重复工作，不能改变答案集合。

\textbf{优化条件。} 本题的关键观察是：单调栈维护当前结果，若栈顶更大且后面还会出现，就可以弹出再等后面放回。这句话必须能翻译成可维护状态；如果题目条件改变后观察不再成立，就不能继续套原来的写法。

\textbf{相邻状态。} 规律是单调栈维护结果字典序，弹出条件必须同时满足栈顶更大且后面还会出现。把这个特例继续拆开看：先构造一个非贪心选择，再比较两者留下的资源、右边界或剩余时间。只有能把非贪心第一步交换成贪心第一步且不变差，局部选择才是规律。

\textbf{状态复用。} 把指数级选择压成一个可证明安全的局部选择；状态通常是当前最远边界、最早结束时间、最小代价或剩余资源。本题落点是：规律是单调栈维护结果字典序，弹出条件必须同时满足栈顶更大且后面还会出现。手算时只改被新输入影响的那一小部分，而不是回头重算完整候选。

\textbf{指针和字段。} 状态字段要能说清：处理顺序、当前已选方案、剩余资源或最远边界、答案计数；若使用排序，要说明排序键；每次局部选择后要能说明当前状态不落后。手算时按这个协议跟踪：拿 cbacdcbc 手算，遇到 b 时，前面的 c 后面还会出现且 c>b，所以可以弹出 c；但遇到只剩最后一次的字符就不能弹。然后按协议复查：手算时同时写一个非贪心选择，与贪心选择比较剩余资源。若无法说明贪心后剩余资源不差，就不能直接下结论。

\textbf{代码落点。} 把局部规则写成业务含义；若需要排序，就处理相同关键字；循环中只维护证明需要的状态，不把局部选择和答案更新混在一起。关键代码行必须能逐句翻译成“旧状态怎样变成新状态”，否则就只是模板。

\textbf{正确性。} 证明从交换或领先性开始：任意最优解都能替换成当前贪心选择而不变差，或贪心状态始终不落后。

\textbf{边界实验。} 先用字符只出现一次、重复字符、弹出后 visited 恢复、字典序比较做反例检查。实验时除了字符只出现一次、重复字符、弹出后 visited 恢复、字典序比较，还要故意替换成一个看似合理但错误的局部规则，构造反例。能解释错误贪心为何失败，才算理解正确贪心。

\textbf{复杂度变化。} 暴力枚举选择顺序可能指数级；贪心通常先排序，再线性扫描或用堆维护少量候选。

\textbf{迁移追问。} 如果题目改成“402 移掉 K 位数字也是单调栈删除，但删除预算和必须保留集合不同”，先判断原状态是否仍然覆盖新答案集合，再决定是微调边界、改 value 语义，还是换算法。

\noindent\textbf{LeetCode 406 根据身高重建队列。}

\textbf{题目说明。} LeetCode 406 根据身高重建队列 的题面入口要先还原成具体任务：高个子的位置先确定后，不会被矮个子影响其前方高个数量。

\textbf{样例入口。} 拿 [7, 0]、[7, 1]、[5, 0] 手算，先放高个子，矮个子插入不会改变高个子前面高个数量；[5, 0] 插到最前也不影响 [7, 0] 的统计。

\textbf{暴力基线。} 慢版本枚举选择顺序或用 DP 保留多个可能方案，先确认最优值存在。随后才检查局部选择是否能通过交换或领先性固定下来。放到本题就是：先按“高个子的位置先确定后，不会被矮个子影响其前方高个数量”完整试慢版本；样例里已经能看到“规律是按身高降序、k 升序排序，再按 k 插入”，所以优化前必须先能说清慢版本枚举了哪些候选。

\textbf{暴力过程。} 拿 [7, 0]、[7, 1]、[5, 0] 手算，先放高个子，矮个子插入不会改变高个子前面高个数量；[5, 0] 插到最前也不影响 [7, 0] 的统计。

\textbf{重复工作。} 题面模型：高个子的位置先确定后，不会被矮个子影响其前方高个数量。暴力版的慢点要落到本题：慢版本会围绕“高个子的位置先确定后，不会被矮个子影响其前方高个数量”枚举选择顺序。样例规律是“规律是按身高降序、k 升序排序，再按 k 插入”，说明存在一个局部选择或边界状态可以安全保留；不能证明这个选择不变差时，就不能把它叫贪心。后面的优化只消掉这类重复工作，不能改变答案集合。

\textbf{优化条件。} 本题的关键观察是：按身高降序、k 升序排序，再按 k 插入当前结果位置。这句话必须能翻译成可维护状态；如果题目条件改变后观察不再成立，就不能继续套原来的写法。

\textbf{相邻状态。} 规律是按身高降序、k 升序排序，再按 k 插入。把这个特例继续拆开看：先构造一个非贪心选择，再比较两者留下的资源、右边界或剩余时间。只有能把非贪心第一步交换成贪心第一步且不变差，局部选择才是规律。

\textbf{状态复用。} 把指数级选择压成一个可证明安全的局部选择；状态通常是当前最远边界、最早结束时间、最小代价或剩余资源。本题落点是：规律是按身高降序、k 升序排序，再按 k 插入。手算时只改被新输入影响的那一小部分，而不是回头重算完整候选。

\textbf{指针和字段。} 状态字段要能说清：处理顺序、当前已选方案、剩余资源或最远边界、答案计数；若使用排序，要说明排序键；每次局部选择后要能说明当前状态不落后。手算时按这个协议跟踪：拿 [7, 0]、[7, 1]、[5, 0] 手算，先放高个子，矮个子插入不会改变高个子前面高个数量；[5, 0] 插到最前也不影响 [7, 0] 的统计。然后按协议复查：手算时同时写一个非贪心选择，与贪心选择比较剩余资源。若无法说明贪心后剩余资源不差，就不能直接下结论。

\textbf{代码落点。} 把局部规则写成业务含义；若需要排序，就处理相同关键字；循环中只维护证明需要的状态，不把局部选择和答案更新混在一起。关键代码行必须能逐句翻译成“旧状态怎样变成新状态”，否则就只是模板。

\textbf{正确性。} 证明从交换或领先性开始：任意最优解都能替换成当前贪心选择而不变差，或贪心状态始终不落后。

\textbf{边界实验。} 先用相同身高、k 为零、插入到末尾、结果稳定性做反例检查。实验时除了相同身高、k 为零、插入到末尾、结果稳定性，还要故意替换成一个看似合理但错误的局部规则，构造反例。能解释错误贪心为何失败，才算理解正确贪心。

\textbf{复杂度变化。} 暴力枚举选择顺序可能指数级；贪心通常先排序，再线性扫描或用堆维护少量候选。

\textbf{迁移追问。} 如果题目改成“这是排序后插入的构造型贪心，证明依赖高个子先固定”，先判断原状态是否仍然覆盖新答案集合，再决定是微调边界、改 value 语义，还是换算法。

\noindent\textbf{LeetCode 435 无重叠区间。}

\textbf{题目说明。} LeetCode 435 无重叠区间 的题面入口要先还原成具体任务：保留最多不重叠区间等价于删除最少区间。

\textbf{样例入口。} 拿 [[1, 2], [2, 3], [3, 4], [1, 3]] 手算，选 [1, 2] 后还能接 [2, 3] 和 [3, 4]；若先选 [1, 3]，会挡住更多区间。

\textbf{暴力基线。} 暴力可以枚举保留哪些区间，再检查它们是否互不重叠，最后用总数减最大保留数量。

\textbf{暴力过程。} 以 [[1, 2], [2, 3], [3, 4], [1, 3]] 为例，选 [1, 2] 后还能接 [2, 3] 和 [3, 4]；若先选 [1, 3] 会挡住更多后续区间。

\textbf{重复工作。} 题面模型：保留最多不重叠区间等价于删除最少区间。暴力版的慢点要落到本题：浪费在枚举保留集合。区间选择里结束越早，给后面留下的空间越大，可以作为安全局部选择。后面的优化只消掉这类重复工作，不能改变答案集合。

\textbf{优化条件。} 本题的关键观察是：按终点升序选择，结束越早给后续留下空间越多。这句话必须能翻译成可维护状态；如果题目条件改变后观察不再成立，就不能继续套原来的写法。

\textbf{相邻状态。} 规律是按右端点升序选择，结束越早给未来留下的空间越多，删除数等于总数减保留数。把这个特例继续拆开看：先构造一个非贪心选择，再比较两者留下的资源、右边界或剩余时间。只有能把非贪心第一步交换成贪心第一步且不变差，局部选择才是规律。

\textbf{状态复用。} 按右端点升序排序。维护 last\_end；若当前区间 start>=last\_end，就保留它并更新 last\_end，否则删除它。手算时只改被新输入影响的那一小部分，而不是回头重算完整候选。

\textbf{指针和字段。} 状态字段要能说清：last\_end 是已保留区间的最后右端，kept 是保留数量，removed 是删除数量。手算时按这个协议跟踪：排序后先保留 [1, 2]；[1, 3] 与它重叠删除；[2, 3] 和 [3, 4] 都能接上。

\textbf{代码落点。} 关键是端点相接是否算重叠。本题中 [1, 2] 和 [2, 3] 不重叠，所以条件是 start>=last\_end。关键代码行必须能逐句翻译成“旧状态怎样变成新状态”，否则就只是模板。

\textbf{正确性。} 任意最优解的第一个区间都可替换为结束最早的区间，后续可用空间不变小；重复此过程得到贪心解。

\textbf{边界实验。} 先用端点相接是否重叠、完全覆盖、相同终点、空列表做反例检查。实验时除了端点相接是否重叠、完全覆盖、相同终点、空列表，还要故意替换成一个看似合理但错误的局部规则，构造反例。能解释错误贪心为何失败，才算理解正确贪心。

\textbf{复杂度变化。} 暴力指数级；排序 O(n log n)，扫描 O(n)，空间取决于排序。

\textbf{迁移追问。} 如果题目改成“用交换论证说明最早结束的选择不劣于其他第一选择”，先判断原状态是否仍然覆盖新答案集合，再决定是微调边界、改 value 语义，还是换算法。

\noindent\textbf{LeetCode 452 用最少数量的箭引爆气球。}

\textbf{题目说明。} LeetCode 452 用最少数量的箭引爆气球 的题面入口要先还原成具体任务：一支箭的位置可以覆盖所有包含该位置的区间。

\textbf{样例入口。} 拿 [[10, 16], [2, 8], [1, 6], [7, 12]] 手算，按右端排序后第一箭射在 6，可以覆盖 [1, 6] 和 [2, 8]；下一箭射在 12。

\textbf{暴力基线。} 暴力可以枚举箭的位置，或对每个气球单独射箭后再尝试合并覆盖关系。

\textbf{暴力过程。} 以 [[10, 16], [2, 8], [1, 6], [7, 12]] 为例，按右端排序后第一支箭射在 6，可同时覆盖 [1, 6] 和 [2, 8]。

\textbf{重复工作。} 题面模型：一支箭的位置可以覆盖所有包含该位置的区间。暴力版的慢点要落到本题：浪费在没有利用区间右端。第一支箭若要覆盖最早结束的气球，射在它右端能给后续重叠区间最大机会。后面的优化只消掉这类重复工作，不能改变答案集合。

\textbf{优化条件。} 本题的关键观察是：按右端点排序，当前箭放在最早结束气球的右端。这句话必须能翻译成可维护状态；如果题目条件改变后观察不再成立，就不能继续套原来的写法。

\textbf{相邻状态。} 规律是每次把箭放在最早结束气球的右端，能最大化后续空间。把这个特例继续拆开看：先构造一个非贪心选择，再比较两者留下的资源、右边界或剩余时间。只有能把非贪心第一步交换成贪心第一步且不变差，局部选择才是规律。

\textbf{状态复用。} 按 right 升序排序。arrow\_pos 初始为第一个区间右端；后续区间若 start<=arrow\_pos 就被当前箭覆盖，否则新增一支箭。手算时只改被新输入影响的那一小部分，而不是回头重算完整候选。

\textbf{指针和字段。} 状态字段要能说清：arrow\_pos 是当前箭的位置，arrows 是箭数，interval 是当前气球闭区间。手算时按这个协议跟踪：第一箭在 6，覆盖 [1, 6] 和 [2, 8]；遇到 [7, 12] 时 7>6，需要第二箭放到 12。

\textbf{代码落点。} 关键是闭区间边界用 start<=arrow\_pos 表示能覆盖。坐标可能为负，不影响排序和比较。关键代码行必须能逐句翻译成“旧状态怎样变成新状态”，否则就只是模板。

\textbf{正确性。} 最早结束区间必须被某支箭覆盖；把这支箭移动到它的右端不会减少已覆盖区间，还可能覆盖更多后续区间。

\textbf{边界实验。} 先用端点相同、负坐标、区间包含、闭区间边界做反例检查。实验时除了端点相同、负坐标、区间包含、闭区间边界，还要故意替换成一个看似合理但错误的局部规则，构造反例。能解释错误贪心为何失败，才算理解正确贪心。

\textbf{复杂度变化。} 暴力选择箭位置复杂；贪心排序 O(n log n)，扫描 O(n)，空间取决于排序。

\textbf{迁移追问。} 如果题目改成“它和无重叠区间是同一类最早结束贪心”，先判断原状态是否仍然覆盖新答案集合，再决定是微调边界、改 value 语义，还是换算法。

\noindent\textbf{LeetCode 630 课程表 III。}

\textbf{题目说明。} LeetCode 630 课程表 III 的题面入口要先还原成具体任务：每门课有持续时间和截止日，已选课程总时长不能超过当前截止日。

\textbf{样例入口。} 拿课程 [100, 200]、[200, 1300]、[1000, 1250] 手算，按截止日排序后若总时长超出当前截止日，就丢掉已选课程里耗时最长的。

\textbf{暴力基线。} 暴力枚举课程子集和排列，检查是否能在各自截止日前完成，取最多课程数。

\textbf{暴力过程。} 课程 [100, 200]、[200, 1300]、[1000, 1250] 中，按截止日处理时，若加入一门课后总时长超出当前截止日，就要考虑删掉已选中耗时最长的课。

\textbf{重复工作。} 题面模型：每门课有持续时间和截止日，已选课程总时长不能超过当前截止日。暴力版的慢点要落到本题：浪费在枚举所有选择。按截止日排序后，当前前缀课程只需要让总耗时尽量小；超时删除最长课程给未来留下最多时间。后面的优化只消掉这类重复工作，不能改变答案集合。

\textbf{优化条件。} 本题的关键观察是：按截止日排序，先选当前课；若超时，就用大顶堆丢掉已选中耗时最长的课。这句话必须能翻译成可维护状态；如果题目条件改变后观察不再成立，就不能继续套原来的写法。

\textbf{相邻状态。} 规律是大顶堆保存已选课程时长，替换最长课能给未来留下最多时间。把这个特例继续拆开看：先构造一个非贪心选择，再比较两者留下的资源、右边界或剩余时间。只有能把非贪心第一步交换成贪心第一步且不变差，局部选择才是规律。

\textbf{状态复用。} 按 deadline 升序排序。扫描课程时先加入 duration 到 maxHeap 并累加 total\_time；若 total\_time>deadline，就弹出堆中最长 duration。手算时只改被新输入影响的那一小部分，而不是回头重算完整候选。

\textbf{指针和字段。} 状态字段要能说清：total\_time 是已选课程总耗时，maxHeap 保存已选课程时长，deadline 是当前必须满足的最紧约束。手算时按这个协议跟踪：加入 1000 后若超过截止日，从已选课程中丢掉耗时最长的一门，课程数量尽量保留，且总耗时降得最多。

\textbf{代码落点。} 关键是大顶堆保存 duration，不保存收益。每次超时只弹一次最长课，因为刚加入前状态已经满足之前截止日。关键代码行必须能逐句翻译成“旧状态怎样变成新状态”，否则就只是模板。

\textbf{正确性。} 在相同课程数量下，总耗时越短越不妨碍后续课程；超时后删除最长课程能得到同数量减一情况下最短总耗时。

\textbf{边界实验。} 先用课程无法单独完成、截止日相同、替换已选课程、空列表做反例检查。实验时除了课程无法单独完成、截止日相同、替换已选课程、空列表，还要故意替换成一个看似合理但错误的局部规则，构造反例。能解释错误贪心为何失败，才算理解正确贪心。

\textbf{复杂度变化。} 暴力指数级；排序 O(n log n)，每门课入堆一次、可能出堆一次，O(n log n) 时间，O(n) 空间。

\textbf{迁移追问。} 如果题目改成“这是反悔贪心：接受后在约束被破坏时删最差候选”，先判断原状态是否仍然覆盖新答案集合，再决定是微调边界、改 value 语义，还是换算法。

\noindent\textbf{LeetCode 646 最长数对链。}

\textbf{题目说明。} LeetCode 646 最长数对链 的题面入口要先还原成具体任务：每个数对像一个区间，选择链要求前一个右端小于后一个左端。

\textbf{样例入口。} 拿 [[1, 2], [2, 3], [3, 4]] 手算，[1, 2] 不能接 [2, 3]，因为要求前一个右端小于后一个左端；可以接 [3, 4]。

\textbf{暴力基线。} 慢版本枚举选择顺序或用 DP 保留多个可能方案，先确认最优值存在。随后才检查局部选择是否能通过交换或领先性固定下来。放到本题就是：先按“每个数对像一个区间，选择链要求前一个右端小于后一个左端”完整试慢版本；样例里已经能看到“规律是按右端升序贪心选择，右端越小越容易接后续数对”，所以优化前必须先能说清慢版本枚举了哪些候选。

\textbf{暴力过程。} 拿 [[1, 2], [2, 3], [3, 4]] 手算，[1, 2] 不能接 [2, 3]，因为要求前一个右端小于后一个左端；可以接 [3, 4]。

\textbf{重复工作。} 题面模型：每个数对像一个区间，选择链要求前一个右端小于后一个左端。暴力版的慢点要落到本题：慢版本会围绕“每个数对像一个区间，选择链要求前一个右端小于后一个左端”枚举选择顺序。样例规律是“规律是按右端升序贪心选择，右端越小越容易接后续数对”，说明存在一个局部选择或边界状态可以安全保留；不能证明这个选择不变差时，就不能把它叫贪心。后面的优化只消掉这类重复工作，不能改变答案集合。

\textbf{优化条件。} 本题的关键观察是：按右端升序选择，当前右端越小越容易接后续数对。这句话必须能翻译成可维护状态；如果题目条件改变后观察不再成立，就不能继续套原来的写法。

\textbf{相邻状态。} 规律是按右端升序贪心选择，右端越小越容易接后续数对。把这个特例继续拆开看：先构造一个非贪心选择，再比较两者留下的资源、右边界或剩余时间。只有能把非贪心第一步交换成贪心第一步且不变差，局部选择才是规律。

\textbf{状态复用。} 把指数级选择压成一个可证明安全的局部选择；状态通常是当前最远边界、最早结束时间、最小代价或剩余资源。本题落点是：规律是按右端升序贪心选择，右端越小越容易接后续数对。手算时只改被新输入影响的那一小部分，而不是回头重算完整候选。

\textbf{指针和字段。} 状态字段要能说清：处理顺序、当前已选方案、剩余资源或最远边界、答案计数；若使用排序，要说明排序键；每次局部选择后要能说明当前状态不落后。手算时按这个协议跟踪：拿 [[1, 2], [2, 3], [3, 4]] 手算，[1, 2] 不能接 [2, 3]，因为要求前一个右端小于后一个左端；可以接 [3, 4]。然后按协议复查：手算时同时写一个非贪心选择，与贪心选择比较剩余资源。若无法说明贪心后剩余资源不差，就不能直接下结论。

\textbf{代码落点。} 把局部规则写成业务含义；若需要排序，就处理相同关键字；循环中只维护证明需要的状态，不把局部选择和答案更新混在一起。关键代码行必须能逐句翻译成“旧状态怎样变成新状态”，否则就只是模板。

\textbf{正确性。} 证明从交换或领先性开始：任意最优解都能替换成当前贪心选择而不变差，或贪心状态始终不落后。

\textbf{边界实验。} 先用端点相等不允许连接、完全覆盖、负数端点、空数组做反例检查。实验时除了端点相等不允许连接、完全覆盖、负数端点、空数组，还要故意替换成一个看似合理但错误的局部规则，构造反例。能解释错误贪心为何失败，才算理解正确贪心。

\textbf{复杂度变化。} 暴力枚举选择顺序可能指数级；贪心通常先排序，再线性扫描或用堆维护少量候选。

\textbf{迁移追问。} 如果题目改成“也可以 DP，但贪心用最早结束证明更简洁”，先判断原状态是否仍然覆盖新答案集合，再决定是微调边界、改 value 语义，还是换算法。

\noindent\textbf{LeetCode 678 有效的括号字符串。}

\textbf{题目说明。} LeetCode 678 有效的括号字符串 的题面入口要先还原成具体任务：星号可以作为左括号、右括号或空，使未匹配左括号数量变成一个可能区间。

\textbf{样例入口。} 拿 (*) 手算，星号可以当空或左括号；拿 (*))，星号必须当左括号才能匹配后面的右括号。与其枚举星号三种选择，不如维护可能未匹配左括号数的区间 [low, high]。

\textbf{暴力基线。} 慢版本枚举选择顺序或用 DP 保留多个可能方案，先确认最优值存在。随后才检查局部选择是否能通过交换或领先性固定下来。放到本题就是：先按“星号可以作为左括号、右括号或空，使未匹配左括号数量变成一个可能区间”完整试慢版本；样例里已经能看到“规律是左括号让区间加一，右括号让区间减一，星号让下界减一上界加一，且下界不能低于零”，所以优化前必须先能说清慢版本枚举了哪些候选。

\textbf{暴力过程。} 拿 (*) 手算，星号可以当空或左括号；拿 (*))，星号必须当左括号才能匹配后面的右括号。与其枚举星号三种选择，不如维护可能未匹配左括号数的区间 [low, high]。

\textbf{重复工作。} 题面模型：星号可以作为左括号、右括号或空，使未匹配左括号数量变成一个可能区间。暴力版的慢点要落到本题：慢版本会围绕“星号可以作为左括号、右括号或空，使未匹配左括号数量变成一个可能区间”枚举选择顺序。样例规律是“规律是左括号让区间加一，右括号让区间减一，星号让下界减一上界加一，且下界不能低于零”，说明存在一个局部选择或边界状态可以安全保留；不能证明这个选择不变差时，就不能把它叫贪心。后面的优化只消掉这类重复工作，不能改变答案集合。

\textbf{优化条件。} 本题的关键观察是：贪心维护可能未匹配左括号数的下界和上界，遇到星号扩大区间并把下界截到零。这句话必须能翻译成可维护状态；如果题目条件改变后观察不再成立，就不能继续套原来的写法。

\textbf{相邻状态。} 规律是左括号让区间加一，右括号让区间减一，星号让下界减一上界加一，且下界不能低于零。把这个特例继续拆开看：先构造一个非贪心选择，再比较两者留下的资源、右边界或剩余时间。只有能把非贪心第一步交换成贪心第一步且不变差，局部选择才是规律。

\textbf{状态复用。} 把指数级选择压成一个可证明安全的局部选择；状态通常是当前最远边界、最早结束时间、最小代价或剩余资源。本题落点是：规律是左括号让区间加一，右括号让区间减一，星号让下界减一上界加一，且下界不能低于零。手算时只改被新输入影响的那一小部分，而不是回头重算完整候选。

\textbf{指针和字段。} 状态字段要能说清：处理顺序、当前已选方案、剩余资源或最远边界、答案计数；若使用排序，要说明排序键；每次局部选择后要能说明当前状态不落后。手算时按这个协议跟踪：拿 (*) 手算，星号可以当空或左括号；拿 (*))，星号必须当左括号才能匹配后面的右括号。与其枚举星号三种选择，不如维护可能未匹配左括号数的区间 [low, high]。然后按协议复查：手算时同时写一个非贪心选择，与贪心选择比较剩余资源。若无法说明贪心后剩余资源不差，就不能直接下结论。

\textbf{代码落点。} 把局部规则写成业务含义；若需要排序，就处理相同关键字；循环中只维护证明需要的状态，不把局部选择和答案更新混在一起。关键代码行必须能逐句翻译成“旧状态怎样变成新状态”，否则就只是模板。

\textbf{正确性。} 证明从交换或领先性开始：任意最优解都能替换成当前贪心选择而不变差，或贪心状态始终不落后。

\textbf{边界实验。} 先用上界变负、最终下界不为零、连续星号、前缀右括号过多做反例检查。实验时除了上界变负、最终下界不为零、连续星号、前缀右括号过多，还要故意替换成一个看似合理但错误的局部规则，构造反例。能解释错误贪心为何失败，才算理解正确贪心。

\textbf{复杂度变化。} 暴力枚举选择顺序可能指数级；贪心通常先排序，再线性扫描或用堆维护少量候选。

\textbf{迁移追问。} 如果题目改成“也可用双栈保存左括号和星号位置，但区间贪心更能说明状态压缩”，先判断原状态是否仍然覆盖新答案集合，再决定是微调边界、改 value 语义，还是换算法。

\noindent\textbf{LeetCode 763 划分字母区间。}

\textbf{题目说明。} LeetCode 763 划分字母区间 的题面入口要先还原成具体任务：每个片段必须覆盖其中所有字符的最后出现位置。

\textbf{样例入口。} 拿 ababcbacadefegdehijhklij 手算，字符 a 的最后位置在 8，所以第一个片段至少到 8；扫描中 b、c 的最后位置也不超过 8，到 8 才能切。

\textbf{暴力基线。} 暴力尝试每个切分位置，并检查切出来的片段中是否有字符还会在后面再次出现。

\textbf{暴力过程。} 字符串 ababcbacadefegdehijhklij 中，a 最后出现在 8，所以第一个片段不可能早于 8 结束；扫描到 b、c 时也要把片段右端扩到它们的最后位置。

\textbf{重复工作。} 题面模型：每个片段必须覆盖其中所有字符的最后出现位置。暴力版的慢点要落到本题：浪费在每次试切后重新检查片段字符。每个字符的最后出现位置可以预处理，扫描时维护当前片段必须覆盖的最远右端。后面的优化只消掉这类重复工作，不能改变答案集合。

\textbf{优化条件。} 本题的关键观察是：扫描维护当前片段最远右端，到达右端即可切分。这句话必须能翻译成可维护状态；如果题目条件改变后观察不再成立，就不能继续套原来的写法。

\textbf{相邻状态。} 规律是当前片段右端是片段内所有字符最后出现位置的最大值。把这个特例继续拆开看：先构造一个非贪心选择，再比较两者留下的资源、右边界或剩余时间。只有能把非贪心第一步交换成贪心第一步且不变差，局部选择才是规律。

\textbf{状态复用。} last[ch] 保存字符最后出现下标。扫描字符串时，right=max(right, last[s[i]])；当 i==right，当前片段可以切开。手算时只改被新输入影响的那一小部分，而不是回头重算完整候选。

\textbf{指针和字段。} 状态字段要能说清：start 是当前片段起点，i 是扫描位置，right 是当前片段内所有字符最后位置的最大值。手算时按这个协议跟踪：第一个片段从 0 开始，right 先因 a 变成 8；扫描到 i=8 时，片段内所有字符都不会再出现，可以切出长度 9。

\textbf{代码落点。} 关键是先预处理 last，再扫描更新 right。切分后记录 right-start+1，并令 start=i+1。关键代码行必须能逐句翻译成“旧状态怎样变成新状态”，否则就只是模板。

\textbf{正确性。} 片段必须覆盖内部每个字符的最后出现位置，所以 right 之前不能切；当 i==right 时，片段内字符都不会越界，立即切能最大化片段数量。

\textbf{边界实验。} 先用单字符、所有字符相同、字符多次跨段做反例检查。实验时除了单字符、所有字符相同、字符多次跨段，还要故意替换成一个看似合理但错误的局部规则，构造反例。能解释错误贪心为何失败，才算理解正确贪心。

\textbf{复杂度变化。} 暴力试切最坏 O(\ensuremath{n^{2}})，预处理加扫描 O(n) 时间，字符集空间 O(1) 或 O(字符集)。

\textbf{迁移追问。} 如果题目改成“这是最早合法切分贪心，能最大化片段数量”，先判断原状态是否仍然覆盖新答案集合，再决定是微调边界、改 value 语义，还是换算法。

\noindent\textbf{LeetCode 860 柠檬水找零。}

\textbf{题目说明。} LeetCode 860 柠檬水找零 的题面入口要先还原成具体任务：找零资源只有 5 元和 10 元，5 元比 10 元更灵活。

\textbf{样例入口。} 拿 bills=[5, 5, 5, 10, 20] 手算，前三个人给 5 后手里有三个 5；收到 10 找一个 5；收到 20 时优先找 10+5，保留更多 5 给后续 10 找零。

\textbf{暴力基线。} 慢版本枚举选择顺序或用 DP 保留多个可能方案，先确认最优值存在。随后才检查局部选择是否能通过交换或领先性固定下来。放到本题就是：先按“找零资源只有 5 元和 10 元，5 元比 10 元更灵活”完整试慢版本；样例里已经能看到“规律是维护 5 元和 10 元数量，20 找零时优先使用较不灵活的 10 元”，所以优化前必须先能说清慢版本枚举了哪些候选。

\textbf{暴力过程。} 拿 bills=[5, 5, 5, 10, 20] 手算，前三个人给 5 后手里有三个 5；收到 10 找一个 5；收到 20 时优先找 10+5，保留更多 5 给后续 10 找零。

\textbf{重复工作。} 题面模型：找零资源只有 5 元和 10 元，5 元比 10 元更灵活。暴力版的慢点要落到本题：慢版本会围绕“找零资源只有 5 元和 10 元，5 元比 10 元更灵活”枚举选择顺序。样例规律是“规律是维护 5 元和 10 元数量，20 找零时优先使用较不灵活的 10 元”，说明存在一个局部选择或边界状态可以安全保留；不能证明这个选择不变差时，就不能把它叫贪心。后面的优化只消掉这类重复工作，不能改变答案集合。

\textbf{优化条件。} 本题的关键观察是：收到 10 用一个 5 找零；收到 20 优先用 10+5，不能时再用三个 5。这句话必须能翻译成可维护状态；如果题目条件改变后观察不再成立，就不能继续套原来的写法。

\textbf{相邻状态。} 规律是维护 5 元和 10 元数量，20 找零时优先使用较不灵活的 10 元。把这个特例继续拆开看：先构造一个非贪心选择，再比较两者留下的资源、右边界或剩余时间。只有能把非贪心第一步交换成贪心第一步且不变差，局部选择才是规律。

\textbf{状态复用。} 把指数级选择压成一个可证明安全的局部选择；状态通常是当前最远边界、最早结束时间、最小代价或剩余资源。本题落点是：规律是维护 5 元和 10 元数量，20 找零时优先使用较不灵活的 10 元。手算时只改被新输入影响的那一小部分，而不是回头重算完整候选。

\textbf{指针和字段。} 状态字段要能说清：处理顺序、当前已选方案、剩余资源或最远边界、答案计数；若使用排序，要说明排序键；每次局部选择后要能说明当前状态不落后。手算时按这个协议跟踪：拿 bills=[5, 5, 5, 10, 20] 手算，前三个人给 5 后手里有三个 5；收到 10 找一个 5；收到 20 时优先找 10+5，保留更多 5 给后续 10 找零。然后按协议复查：手算时同时写一个非贪心选择，与贪心选择比较剩余资源。若无法说明贪心后剩余资源不差，就不能直接下结论。

\textbf{代码落点。} 把局部规则写成业务含义；若需要排序，就处理相同关键字；循环中只维护证明需要的状态，不把局部选择和答案更新混在一起。关键代码行必须能逐句翻译成“旧状态怎样变成新状态”，否则就只是模板。

\textbf{正确性。} 证明从交换或领先性开始：任意最优解都能替换成当前贪心选择而不变差，或贪心状态始终不落后。

\textbf{边界实验。} 先用第一张不是 5、连续 10、连续 20、只有 5、5 元数量不足做反例检查。实验时除了第一张不是 5、连续 10、连续 20、只有 5、5 元数量不足，还要故意替换成一个看似合理但错误的局部规则，构造反例。能解释错误贪心为何失败，才算理解正确贪心。

\textbf{复杂度变化。} 暴力枚举选择顺序可能指数级；贪心通常先排序，再线性扫描或用堆维护少量候选。

\textbf{迁移追问。} 如果题目改成“若面额集合变化，需要重新证明哪种资源更灵活”，先判断原状态是否仍然覆盖新答案集合，再决定是微调边界、改 value 语义，还是换算法。

\subsection{实验复盘}

追问通常会要求证明。请准备交换论证或领先性论证，并给出一个看似合理但错误的贪心规则反例。

复盘本节时先从 LeetCode 45 跳跃游戏 II、LeetCode 55 跳跃游戏、LeetCode 122 买卖股票 II 开始：每道题都先手写一个能覆盖全部候选的暴力版，再指出优化结构具体保存了哪一段历史信息，最后用相邻案例比较为什么不能互换解法。
